% 11-711 Fall 2026; Homework 1 report template.
%
% Fill in each section. Written questions are graded ONLY from this document.
% The autograder never sees it, and we do not go looking in your code for
% answers. Answer every part; if you skipped a question, say so rather than
% leaving it blank, so we can tell the difference between skipped and missing.
%
% Per the syllabus, every word here must be your own. You may use AI to help
% write code, but not to write any part of this report.

\documentclass[11pt]{article}

\usepackage[margin=1in]{geometry}
\usepackage{xcolor}
\usepackage{titlesec}
\usepackage{graphicx}
\usepackage{amsmath,amssymb}
\usepackage{listings}
\usepackage{lmodern}

\definecolor{mauve}{rgb}{0.58,0,0.82}
\titleformat{\section}[hang]{\color{mauve}\normalfont\sffamily\Large}{\thesection}{1em}{}
\titleformat{\subsection}[hang]{\color{mauve}\normalfont\sffamily\large}{\thesubsection}{1em}{}

\lstset{basicstyle=\ttfamily\small, breaklines=true, frame=single,
        columns=flexible, showstringspaces=false}

\title{11-711 Homework 1}
\author{Your Name \\ \texttt{yourandrewid}}
\date{Due: September 24, 2026}

% Long \texttt{} paths cannot hyphenate, which pushes lines past the right margin.
% emergencystretch lets TeX stretch interword space as a last resort rather than
% overflowing; \sloppy prefers loose lines to overfull ones.
\emergencystretch=3em
\sloppy
\hbadness=10000

\begin{document}
\maketitle

%======================================================================
\section*{Problem 1: Tokenization}

\subsection*{Question 1.1 (12 points)}
\subsubsection*{1.1.A, counterexamples}
% At most two sentences per property (injective, invertible,
% concatenation-preserving) describing the violation you found with
% google-bert/bert-base-cased.

\subsubsection*{1.1.B, why concatenation cannot be preserved}

\subsection*{Question 1.3 (4 points)}
\subsubsection*{1.3.A, the longest token}
\subsubsection*{1.3.B, BPE and training-data privacy}

\subsection*{Question 1.4 (4 points)}
\subsubsection*{1.4.A, digits}
\subsubsection*{1.4.B, domains}
% A small table of tokens-per-character is the clearest way to present this.
\subsubsection*{1.4.C, who pays}

%======================================================================
\section*{Problem 2: Implementing Llama-2}

\subsection*{Question 2.2 (6 points)}
\subsubsection*{2.2.A, what weight tying does}
\subsubsection*{2.2.B, one line of PyTorch}
\begin{lstlisting}[language=Python]
# your one line here
\end{lstlisting}

\subsection*{Question 2.3 (7 points)}
\subsubsection*{2.3.A, KV cache size}
\subsubsection*{2.3.B, the two numbers}
\subsubsection*{2.3.C, what GQA buys, and its cost}

%======================================================================
\section*{Problem 3: Pretraining}

\subsection*{Question 3.1 (5 points)}
\subsubsection*{3.1.A, why $\log(8)$}
\subsubsection*{3.1.B, a bug that would beat the floor}
\subsubsection*{3.1.C, why 3.4–3.6 would be meaningless}

\subsection*{Question 3.3 (6 points)}
\subsubsection*{3.3.A, the two expressions}
\subsubsection*{3.3.B, why warm up, why anneal}

\subsection*{Question 3.4 (5 points)}
\subsubsection*{3.4.A, final validation loss}
\subsubsection*{3.4.B, steps to saturation}
\subsubsection*{3.4.C, training curve}
% \includegraphics[width=0.7\textwidth]{llama-tiny-curve.png}

\subsection*{Question 3.5 (14 points)}
% At most THREE configurations. A table is the clearest format:
%
%   Config | what changed | FLOPs | val perplexity
%
\subsubsection*{3.5.A, experimental procedure}
\subsubsection*{3.5.B, results}
\subsubsection*{3.5.C, your best configuration}
\begin{lstlisting}
# paste your YAML here
\end{lstlisting}
\subsubsection*{3.5.D, size versus steps at fixed compute}

\subsection*{Question 3.6 (reported here, graded on code)}
% Final validation loss and perplexity.

%======================================================================
\section*{Problem 4: Generation and perplexity}

\subsection*{Question 4.2 (11 points)}
\subsubsection*{4.2.A, prefixes, settings, and generations}
\subsubsection*{4.2.B, the effect of temperature}
\subsubsection*{4.2.C, your model versus stories42M}
\subsubsection*{4.2.D, what you cannot conclude}

\subsection*{Question 4.3 (12 points)}
\subsubsection*{4.3.A, removing ``the''}
\subsubsection*{4.3.B, distribution shift}

%======================================================================
\section*{Problem 5: Adapting the model}

\subsection*{Question 5.2 (8 points)}
\subsubsection*{5.2.A, dev accuracy, four runs}
% Report dev only. cfimdb-test.txt ships with placeholder labels.
\subsubsection*{5.2.B, majority baselines, prediction distributions, diagnosis}
\subsubsection*{5.2.C, what finetuning changed}

\subsection*{Question 5.3 (7 points)}
\subsubsection*{5.3.A, your model versus stories42M}
\subsubsection*{5.3.B, difference versus run-to-run variation}
\subsubsection*{5.3.C, what pretraining actually provides}

%======================================================================
\section*{Question 6: Use of AI (required, ungraded)}
% Which parts of this assignment did you use AI for, and how? If you used none,
% say so. This question is required by the syllabus: using AI without making an
% honest attempt at answering it is itself an academic integrity violation.

%======================================================================
\section*{Bonus (optional, up to 10 points)}
% Which track, what you did, what you measured, what you concluded.
% Include plots and training logs.

\end{document}
